% Provenance: user-supplied 56-page V5 PDF, Tables 7-11, 2026-11-12.
% Transcribed results, not independently rerun. This supplementary file
% does not replace the native V5 LaTeX author source.
\section{V5 Experimental Results: Preservation and Reconciliation}
\label{sec:v5_exp_preserve}
The numerical results below are preserved from the V5 manuscript, including
the new cache simulations, retention ablation and recovery-logic ablation.
They remain \emph{author-reported}; a complete independent reproduction
requires raw event logs and per-instance predictions. The simulated cache
results come from a CPU model, not a production GPU serving engine.

\subsection{Cross-layer simulation results}
The first scenario used three agents, a 100-block cache, 50 conversations,
seed 42 and a five-conversation warm-up. It explicitly assumes diagnostic
content and serving cache blocks occupy a constrained shared buffer.
\begin{table*}[ht]\centering\footnotesize
\caption{V5 Table 7: author-reported simulated cache policies. Hit-rate
and eviction event streams require code-level reconciliation.}
\label{tab:v5sim1}
\begin{tabular}{lrrrr}\toprule
Policy & Evictions & Diag-bearing & Hit rate & Diag retained\\\midrule
LRU & 2,688 & 58.5\% & 100\% & 67/100\\
FIFO & 2,641 & 60.3\% & 100\% & 64/100\\
Diag-aware & 2,658 & 58.2\% & 88.4\% & 100/100\\
\bottomrule\end{tabular}\end{table*}
The second scenario comprises 20 conversations, 240 requests, a
5,000-token cache, and a separate workload; it must not be compared as
though it used the first scenario's cache budget.
\begin{table*}[ht]\centering\footnotesize
\caption{V5 Table 8: author-reported CPU serving simulation.}
\label{tab:v5sim2}
\begin{tabular}{lrrrr}\toprule
Policy & Prefix hit & Utilization & Mean TTFT (ms) & Evictions\\\midrule
LRU Flat & 95.45\% & 91.1\% & 13.456 & 208\\
RadixTree & 95.74\% & 25.6\% & 13.240 & 0\\
\bottomrule\end{tabular}\end{table*}
The access/hit/miss definitions must be reconciled with the reported 2,688
evictions and 100\% hit rate; cumulative diagnostic-bearing eviction
events must be distinguished from final retained diagnostic blocks.

\subsection{Real-trace retention ablation}
The source reports 220 annotated TraceElephant failure traces:
Captain-Agent/GAIA (73), Captain-Agent/AssistantBench (12),
Magentic-One/GAIA (74), Magentic-One/AssistantBench (17),
SWE-Agent/SWE-Bench Verified (44). The policies are S1 full retention,
S2 random sampling, S3 recency and S4 importance scoring; attribution
uses both causal-graph and last-error heuristics.
\begin{table*}[ht]\centering\footnotesize
\caption{V5 Table 9: 10\% retention accuracy. S2 is mean $\pm$ SD over
three seeds; other methods are deterministic.}
\label{tab:v5ret10}
\begin{tabular}{lrr}\toprule
Policy & Causal graph & Last-error\\\midrule
S1 Full & .459 & .432\\
S2 Random & $.524\pm.023$ & $.523\pm.020$\\
S3 Recency & .391 & .391\\
S4 Importance & .532 & .500\\
\bottomrule\end{tabular}\end{table*}
\begin{table*}[ht]\centering\scriptsize
\caption{V5 Table 10: causal-graph accuracy over the reported
10--100\% retention budgets; author-reported.}
\label{tab:v5retall}
\begin{tabular}{lrrrr}\toprule
Budget & S1 Full & S2 Random & S3 Recency & S4 Importance\\\midrule
10\% & .459 & .524 & .391 & .532\\
20\% & .459 & .479 & .423 & .555\\
30\% & .459 & .479 & .423 & .546\\
40\% & .459 & .479 & .445 & .514\\
50\% & .459 & .461 & .450 & .496\\
60\% & .459 & .450 & .459 & .491\\
70\% & .459 & .432 & .459 & .482\\
80\% & .459 & .451 & .464 & .473\\
90\% & .459 & .457 & .459 & .464\\
100\% & .459 & .459 & .459 & .459\\
\bottomrule\end{tabular}\end{table*}

The inverted curve largely reflects an attributor-specific recovery
heuristic that suppresses earlier failures after later normal messages,
rather than demonstrating that discarding informative observations is
universally beneficial. Importance scoring differs by only +0.008
from random sampling at 10\% budget (within reported noise).

\subsection{Corrected interpretation of paired recovery ablation}
\begin{table*}[ht]\centering\footnotesize
\caption{V5 Table 11: reported accuracies retained, while the
two-sided \emph{exact} McNemar p-value and discordant-pair direction
are reconciled. Reported CIs require independent reproduction from
per-case outputs.}
\label{tab:v5recover}
\begin{tabular}{lrrr}\toprule
Condition & Correct & Accuracy & Exact two-sided $p$\\\midrule
C1 Original, 100\% retention & 101/220 & .459 & --\\
C2 No recovery logic, 100\% & 117/220 & .532 & .003719 vs C1\\
C3 C2 + S4 at 10\% & 119/220 & .541 & .891923 vs C2\\
\bottomrule\end{tabular}\end{table*}
Define $b$ as new-only correct and $c$ as old-only correct:
C1$\to$C2 must have $b=22,c=6$ (net 16);
C2$\to$C3 must have $b=28,c=26$ (net 2).
Those two p-values are two-sided \emph{exact} binomial McNemar values.
If a two-sided mid-$p$ procedure were intended, the corresponding
values are .0023157001 and .7877061897. The extra 0.9 percentage points
from C2 to C3 are not statistically significant; the paper's paired
confidence interval should be recalculated from raw outcome vectors.

\subsection{DRICE synthesis}
The V5 DRICE framework includes five orthogonal dimensions:
D (diagnostic evidence and capture), R (retention trigger and diagnostic
awareness), I (model/workload/memory assumptions), C (compute, memory,
latency and communication cost), and E (source evidence depth).
Claims that a system jointly optimizes diagnostic evidence and serving
state require an explicit shared resource budget and an observed
diagnostic-fidelity metric, not only ordinary prefix-cache hits.
A complete evidence matrix must attach claim-level section/table/page
locators to each DRICE cell. See the GitHub V5_EVIDENCE directory.

\paragraph*{Source fidelity.} These tables preserve V5 experimental
numbers and their interpretation. The V5_SOURCE_SNAPSHOT text files
preserve the manuscript page-by-page; a native editable V5 LaTeX
archive with its current graphics is still required before guaranteeing
full Overleaf manuscript fidelity.
